% ============================================================
% DOCUMENTCLASS
% ============================================================
\documentclass[12pt,oneside]{book}

% ============================================================
% METADATOS
% ============================================================
% NOTA EDITORIAL: el apunte original no indica facultad, carrera ni nombre de la materia.
% NOTA EDITORIAL: la primera parte de la transcripción (conversaciones sobre casos concretos
% de consultantes, con datos personales identificables) no contiene desarrollo académico y se omitió.
\newcommand{\universidad}{Universidad de Buenos Aires (UBA)}
\newcommand{\materia}{Apuntes de estudio}
\newcommand{\titulo}{Proceso civil, mediaci\'on y atenci\'on de consultantes}
\newcommand{\autor}{Isaac Edgar Camacho Ocampo}
\newcommand{\anio}{2026}

% ============================================================
% PAQUETES
% ============================================================
\usepackage[utf8]{inputenc}
\usepackage[T1]{fontenc}
\usepackage[spanish,es-nodecimaldot]{babel}

\usepackage[
    top=2.5cm,
    bottom=2.5cm,
    left=3cm,
    right=2.5cm,
    headheight=15pt
]{geometry}

\usepackage{amsmath}
\usepackage{amssymb}
\usepackage{mathtools}

\usepackage{graphicx}
\usepackage{float}
\usepackage{caption}
\usepackage{subcaption}

\usepackage{booktabs}
\usepackage{array}
\usepackage{longtable}
\usepackage{tabularx}

\usepackage{enumitem}

\usepackage{xcolor}
\usepackage[most]{tcolorbox}

\usepackage{fancyhdr}
\usepackage{ragged2e}

\usepackage[
    colorlinks=true,
    linkcolor=blue!50!black,
    citecolor=green!40!black,
    urlcolor=blue!60!black,
    pdftitle={Proceso civil, mediación y atención de consultantes},
    pdfauthor={Isaac Edgar Camacho Ocampo},
    pdfsubject={Apuntes universitarios},
    bookmarks=true
]{hyperref}

% ============================================================
% CONFIGURACIÓN
% ============================================================
\graphicspath{
    {./img/}
    {./graficos/}
    {./figuras/}
}

\setcounter{tocdepth}{2}

\setlength{\parindent}{0pt}
\setlength{\parskip}{0.5em}

\setlist[itemize]{
    topsep=0.3em,
    itemsep=0.2em
}

\setlist[enumerate]{
    topsep=0.3em,
    itemsep=0.2em
}

\clubpenalty=10000
\widowpenalty=10000

\pagestyle{fancy}
\fancyhf{}
\fancyhead[L]{\nouppercase{\leftmark}}
\fancyhead[R]{\materia}
\fancyfoot[C]{\thepage}
\renewcommand{\headrulewidth}{0.4pt}
\renewcommand{\footrulewidth}{0pt}

\fancypagestyle{plain}{
    \fancyhf{}
    \fancyfoot[C]{\thepage}
    \renewcommand{\headrulewidth}{0pt}
}

% ============================================================
% COMANDOS
% ============================================================
\newcommand{\abs}[1]{\left|#1\right|}
\newcommand{\norm}[1]{\left\lVert#1\right\rVert}
\newcommand{\vect}[1]{\mathbf{#1}}
\newcommand{\deriv}[2]{\frac{d#1}{d#2}}
\newcommand{\pderiv}[2]{\frac{\partial#1}{\partial#2}}

% ============================================================
% ENTORNOS / CAJAS
% ============================================================
\newtcolorbox{definition}[1]{
    enhanced,
    breakable,
    title={Definición: #1},
    fonttitle=\bfseries,
    colback=blue!3,
    colframe=blue!50!black,
    boxrule=0.8pt,
    arc=2mm
}

\newtcolorbox{important}{
    enhanced,
    breakable,
    title={Importante},
    fonttitle=\bfseries,
    colback=yellow!8,
    colframe=orange!70!black,
    boxrule=0.8pt,
    arc=2mm
}

\newtcolorbox{example}[1]{
    enhanced,
    breakable,
    title={Ejemplo: #1},
    fonttitle=\bfseries,
    colback=green!4,
    colframe=green!40!black,
    boxrule=0.8pt,
    arc=2mm
}

\newtcolorbox{formula}[1]{
    enhanced,
    breakable,
    title={Fórmula / Regla: #1},
    fonttitle=\bfseries,
    colback=purple!4,
    colframe=purple!50!black,
    boxrule=0.8pt,
    arc=2mm
}

\newtcolorbox{exercise}[1]{
    enhanced,
    breakable,
    title={Ejercicio: #1},
    fonttitle=\bfseries,
    colback=gray!5,
    colframe=gray!60!black,
    boxrule=0.8pt,
    arc=2mm
}

\newtcolorbox{note}{
    enhanced,
    breakable,
    title={Nota},
    fonttitle=\bfseries,
    colback=white,
    colframe=gray!50,
    boxrule=0.6pt,
    arc=2mm
}

% ============================================================
% DOCUMENTO
% ============================================================
\begin{document}

% ------------------------------------------------------------
% PORTADA
% ------------------------------------------------------------
\begin{titlepage}

\thispagestyle{empty}

\begin{center}

\vspace*{1cm}

{\large \textsc{\universidad}}

\vspace{3.5cm}

{\Huge \textbf{\titulo}}

\vspace{1cm}

\textit{Comprender los conceptos es el primer paso para convertir el estudio en conocimiento.}

\vspace{1.2cm}

\rule{0.70\textwidth}{0.5pt}

\vspace{0.5cm}

{\large Apuntes de estudio}

\vspace{0.5cm}

\rule{0.70\textwidth}{0.5pt}

\vfill

{\large \autor}

\vspace{0.3cm}

{\large \anio}

\end{center}

\vfill

\begin{flushleft}
\textit{Este es un modesto aporte a los estudiantes y no pretende sustituir el material oficial de la materia.}
\end{flushleft}

\end{titlepage}

% ------------------------------------------------------------
% ÍNDICE
% ------------------------------------------------------------
\tableofcontents

% ------------------------------------------------------------
% CONTENIDO
% ------------------------------------------------------------

% ============================================================
\chapter[El proceso civil y los recursos]{El proceso civil y el sistema de recursos}
% ============================================================

\section{Línea de tiempo del proceso civil}

Para facilitar la comprensión del proceso civil se utiliza una línea de tiempo de elaboración propia de la cátedra. Se aclara que este gráfico no tiene rigor científico, pero ayuda a representar mentalmente el desarrollo del proceso. Sus puntos más importantes son aquellos que deben conocerse al finalizar la formación.

Las etapas que se señalan, en orden, son:

\begin{enumerate}
    \item \textbf{Mediación}. Algunos autores la ubican fuera del proceso civil; en la línea de tiempo utilizada se la incluye como primera etapa.
    \item \textbf{Demanda}.
    \item \textbf{Contestación de demanda}.
    \item \textbf{Apertura a prueba}.
    \item \textbf{Alegatos}.
    \item \textbf{Sentencia}.
\end{enumerate}

\section{Los recursos}

A las etapas anteriores se agregan los \textbf{recursos}, que no se ubican en un punto único de la línea de tiempo sino que abarcan todo el proceso. La razón es que puede revocarse una providencia simple, es decir, se puede cuestionar una resolución de ese tipo en distintos momentos del trámite.

\begin{itemize}
    \item Los recursos operan desde el inicio de la demanda y hasta después de los alegatos.
    \item Contra la sentencia, el recurso mencionado es la \textbf{apelación}, señalada como el único recurso contra una sentencia.
    \item Respecto de la revocatoria de providencias simples, se mencionan las denominaciones \emph{revocatoria} y \emph{reposición}, y se advierte que no son lo mismo y que no importa fijarse solo en los nombres, pues se verán más adelante.
\end{itemize}

% TODO: contenido incompleto o ambiguo en las notas originales (denominaciones "revocatoria", "reposición" y "revocatoria con apelación": la transcripción no permite fijar la distinción exacta)

\begin{table}[H]
\centering
\begin{tabularx}{\textwidth}{>{\raggedright\arraybackslash}p{0.28\textwidth}>{\raggedright\arraybackslash}X}
\toprule
\textbf{Etapa o elemento} & \textbf{Observación según el apunte} \\
\midrule
Mediación & Algunos autores la colocan fuera del proceso civil. \\
Demanda, contestación, apertura a prueba, alegatos & Etapas centrales del proceso civil. \\
Sentencia & Contra ella procede el recurso de apelación. \\
Recursos & Abarcan todo el proceso: desde el inicio de la demanda hasta después de los alegatos (providencias simples). \\
\bottomrule
\end{tabularx}
\caption{Síntesis de la línea de tiempo del proceso civil.}
\end{table}

\section{Cuadro Cornell del capítulo}

\begin{longtable}{
    >{\raggedright\arraybackslash}p{0.26\textwidth}
    >{\raggedright\arraybackslash}p{\dimexpr\textwidth-0.26\textwidth-4\tabcolsep\relax}
}
\toprule
\textbf{Preguntas / palabras clave} & \textbf{Notas / respuestas} \\
\midrule
\endhead

\textbf{¿Cuáles son las etapas que se grafican en la línea de tiempo del proceso civil?} &
Mediación, demanda, contestación de demanda, apertura a prueba, alegatos y sentencia. \\[1.5em]

\textbf{¿Qué discusión existe sobre la mediación?} &
Algunos autores la ubican fuera del proceso civil; en la línea de tiempo de clase se la incluye. \\[1.5em]

\textbf{¿Por qué los recursos se grafican abarcando todo el proceso?} &
Porque puede revocarse una providencia simple, por lo que los recursos operan desde el inicio de la demanda hasta después de los alegatos. \\[1.5em]

\textbf{¿Qué recurso procede contra la sentencia?} &
La apelación, indicada como el único recurso contra una sentencia. \\[1.5em]

\midrule
\multicolumn{2}{p{\dimexpr\textwidth-2\tabcolsep\relax}}{%
\textbf{Resumen:} El proceso civil se grafica con una línea de tiempo (mediación, demanda, contestación, apertura a prueba, alegatos y sentencia) sin rigor científico pero útil para visualizarlo. Los recursos abarcan todo el proceso para las providencias simples, y contra la sentencia procede la apelación.} \\
\bottomrule
\end{longtable}

% ============================================================
\chapter[Atención de consultantes]{Deberes en la atención de consultantes}
% ============================================================

\section{No se garantizan resultados}

En el ámbito del patrocinio no se da ningún tipo de resultado. Jamás se le dice a una persona si va a ganar su juicio; sí se le informa qué posibilidades existen de iniciarlo y de continuarlo.

\begin{important}
El abogado que ofrece resultados está estafando. Aun cuando tenga mucha fama o visibilidad, no puede asegurar un resultado, porque en un juicio intervienen muchos factores, y mucho más en materia civil, patrimonial y de familia. Se señala que el ámbito penal es otra historia. Siempre habrá algún factor que impida garantizar el resultado.
\end{important}

Se plantea, y la docente lo confirma, que la obligación del abogado es una obligación de \textbf{medios} y no de resultados.

Esta pauta constituye una obligación dentro de la clase; luego de recibidos, se la presenta como un consejo para quien decida tomarlo.

\section{Deber de privacidad}

Se puede conversar sobre casos, pero \textbf{no se pueden dar nombre, apellido ni referencias que identifiquen a las personas}. Por ejemplo, si se trata de una persona con una discapacidad severa, no debe mencionarse esa discapacidad si con ello puede saberse de quién se habla.

\begin{itemize}
    \item No hace falta aportar nombre y apellido para comentar un caso; la comparación que se utiliza es la del psicólogo.
    \item Tampoco debe darse la localidad.
    \item Si se desea comentar un caso con terceros fuera del ámbito de la clase, debe hacerse con esta particularidad y con especial cuidado, sobre todo con otros colegas, porque la confidencialidad se compromete fácilmente: no se sabe quién conoce a quién.
    \item Está en juego la responsabilidad y la matrícula de quienes ya son profesionales; en el caso de los estudiantes, la de la cátedra.
    \item Los abogados son muchos pero a la vez pocos, y a veces se encuentran con letrados conocidos en la parte contraria.
\end{itemize}

\section{Comunicación con los consultantes}

\begin{itemize}
    \item No se contestan mails ni mensajes de WhatsApp, salvo que se indique expresamente algo puntual (por ejemplo, avisar determinada cuestión). Los encuentros se realizan en los días de comisión, lunes y jueves.
    \item No se mantienen conversaciones con los consultantes por esas vías. La razón es cuidar a los estudiantes: se desconoce lo que ocurre en esas conversaciones y pueden aparecer personas que no tratan bien o que reaccionan de modo violento.
    \item La conversación debe ser \textbf{estrictamente profesional}.
    \item El procedimiento es: se redacta un escrito, se firma y se trae a la comisión según lo indicado.
    \item Aun cuando el consultante provoque enojo o no quiera atender explicaciones, se le indica que el asunto se deja en manos de la comisión.
\end{itemize}

\section{No se atiende fuera del ámbito de la facultad}

A los consultantes \textbf{no se los atiende fuera} de la facultad, ni siquiera brevemente para recibir una firma, ni con un saludo. Es equivalente a lo indicado para WhatsApp: no se puede saber qué sucede ni entender lo ocurrido.

\begin{itemize}
    \item Si un consultante reconoce a un estudiante fuera de la facultad y se acerca, el estudiante debe indicarle que no se lo atiende afuera y que se lo atenderá en el ámbito de la comisión.
    \item Se mencionan situaciones, no violentas pero sí de confusión, en las que la cátedra tuvo que intervenir.
\end{itemize}

\section{Cuadro Cornell del capítulo}

\begin{longtable}{
    >{\raggedright\arraybackslash}p{0.26\textwidth}
    >{\raggedright\arraybackslash}p{\dimexpr\textwidth-0.26\textwidth-4\tabcolsep\relax}
}
\toprule
\textbf{Preguntas / palabras clave} & \textbf{Notas / respuestas} \\
\midrule
\endhead

\textbf{¿Qué se le puede decir a un consultante sobre el resultado de su juicio?} &
Nunca se garantiza ni se anticipa un resultado; solo se informan las posibilidades de iniciar y continuar el caso. El abogado que ofrece resultados está estafando, y la obligación es de medios. \\[1.5em]

\textbf{¿Por qué no puede asegurarse un resultado?} &
Porque en un juicio, especialmente civil, patrimonial y de familia, intervienen muchos factores. \\[1.5em]

\textbf{¿En qué consiste el deber de privacidad?} &
En no dar nombre, apellido ni referencias que permitan identificar a las personas (por ejemplo, una discapacidad severa o la localidad) al hablar de casos, con especial cuidado frente a otros colegas. \\[1.5em]

\textbf{¿Cómo debe ser la comunicación con los consultantes?} &
Estrictamente profesional: sin responder mails ni WhatsApp salvo indicación puntual; se pasa un escrito, se firma y se trae a la comisión. \\[1.5em]

\textbf{¿Se puede atender a un consultante fuera de la facultad?} &
No: ni siquiera para una firma o un saludo. Si se acerca, se le indica que no se atiende afuera. \\[1.5em]

\midrule
\multicolumn{2}{p{\dimexpr\textwidth-2\tabcolsep\relax}}{%
\textbf{Resumen:} En la atención de consultantes no se garantizan resultados (obligación de medios), se preserva estrictamente la privacidad de los casos, la comunicación es exclusivamente profesional y se realiza en el ámbito de la comisión, nunca por mensajería ni fuera de la facultad.} \\
\bottomrule
\end{longtable}

% ============================================================
\chapter{La mediación}
% ============================================================

\section{Servicio de mediación}

Se informa que existe un servicio de mediación gratuita. El consultante solo abona la carta documento, que se estima en alrededor de 25.000 pesos, de modo que no hay otros gastos para los consultantes. Este servicio se usaba mucho y su uso ha ido disminuyendo.

Actualmente se utiliza, con resultados poco favorables, un centro de justicia de Nación ubicado en la avenida Córdoba, que permite mediar de forma presencial y virtual. Su uso implica un trámite de registro y un formulario cuyo manejo se está aprendiendo.

% TODO: contenido incompleto o ambiguo en las notas originales (monto de la carta documento expresado como estimación)

\section{Obligatoriedad y excepciones}

La mediación es obligatoria cuando la ley lo dispone. Se identifican los artículos 1, 2, 3, 4 y 8 como los más importantes, porque determinan qué se media y qué no.

% TODO: contenido incompleto o ambiguo en las notas originales (el número de la ley se dicta de forma confusa: "26 182", luego "26 58 58"; no se completa con conocimiento externo)

\begin{itemize}
    \item \textbf{Lo que no se puede mediar} se inicia directamente. Se menciona como ejemplo que las medidas cautelares se inician directamente.
    \item \textbf{Para todo lo demás hay que mediar.}
    \item En el fuero de familia se indica que es uno de los fueros en que más se media: alimentos, régimen de comunicación y otras cuestiones.
\end{itemize}

% TODO: contenido incompleto o ambiguo en las notas originales (frase "Alimentos prohi..." sobre medida cautelar; y enumeración "cálculos pro vida" en el fuero de familia)

\begin{example}{Medidas cautelares}
Cuando un edificio está por caerse no se va a juicio sino que se inicia una medida cautelar. Del mismo modo, cuando una persona no tiene cuota alimentaria y no tiene para sobrevivir, se inicia una medida cautelar y se promueve directamente el juicio, sin pasar por mediación.
\end{example}

\section{El formulario de mediación}

La mediación se inicia mediante un formulario con los datos de las partes, disponible en el drive de la comisión.

\begin{itemize}
    \item El formulario debe completarse \textbf{por completo}: los mediadores suelen rebotar los formularios a la comisión, lo que genera inconvenientes frente a los consultantes (por ejemplo, ante un domicilio mal consignado).
    \item No debe dejarse ningún espacio en blanco, porque son precisamente los espacios en blanco los que se controlan. Si no se conoce un dato, se avisa a la docente.
    \item En el campo de mail del demandado, si no se lo conoce, se coloca el de la docente, el de la comisión o se indica que se desconoce.
    \item Si existen \textbf{medidas cautelares} (de violencia, de régimen de comunicación o de alimentos), deben consignarse en el apartado especial, y en particular la de violencia, para que el mediador esté al tanto.
\end{itemize}

\section{Cuadro Cornell del capítulo}

\begin{longtable}{
    >{\raggedright\arraybackslash}p{0.26\textwidth}
    >{\raggedright\arraybackslash}p{\dimexpr\textwidth-0.26\textwidth-4\tabcolsep\relax}
}
\toprule
\textbf{Preguntas / palabras clave} & \textbf{Notas / respuestas} \\
\midrule
\endhead

\textbf{¿Cuánto debe pagar el consultante en la mediación gratuita?} &
Solo la carta documento, estimada en unos 25.000 pesos; no hay otros gastos. \\[1.5em]

\textbf{¿Cuándo es obligatoria la mediación?} &
Cuando la ley lo dispone; los artículos 1, 2, 3, 4 y 8 son los que indican qué se media y qué no. \\[1.5em]

\textbf{¿Qué se inicia directamente sin mediar?} &
Lo que no se puede mediar, por ejemplo las medidas cautelares (como cuando un edificio está por caerse o una persona no tiene cuota alimentaria). Para todo lo demás hay que mediar. \\[1.5em]

\textbf{¿Cómo debe completarse el formulario de mediación?} &
Por completo, sin espacios en blanco; si falta un dato se avisa a la docente; si se desconoce el mail del demandado se coloca el de la comisión o se indica que se desconoce. \\[1.5em]

\textbf{¿Qué debe consignarse respecto de las medidas cautelares?} &
Las medidas de violencia, régimen de comunicación o alimentos deben ponerse en el apartado especial, en particular la de violencia. \\[1.5em]

\midrule
\multicolumn{2}{p{\dimexpr\textwidth-2\tabcolsep\relax}}{%
\textbf{Resumen:} La mediación es obligatoria cuando la ley lo dispone y lo que no puede mediarse (como las medidas cautelares) se inicia directamente. El servicio es gratuito salvo la carta documento, y el formulario debe completarse sin espacios en blanco, informando especialmente las medidas cautelares, sobre todo las de violencia.} \\
\bottomrule
\end{longtable}

% ============================================================
\chapter[Protocolo de violencia]{Protocolo de violencia en la mediación}
% ============================================================

Existe un protocolo de violencia que debe aplicarse cuando llega un formulario en el que se consigna una causa de violencia. Se pide inmediatez: se avisa inmediatamente y los estudiantes actúan sin demora.

\section{Pasos del protocolo}

\begin{enumerate}
    \item \textbf{Llamar de inmediato a la consultante o al consultante.} Se llama de manera cubierta: se hace por una solicitud genérica de la UBA, sin dar muchos datos y sin mencionar el patrocinio gratuito, porque no se sabe si la persona agresora está cerca. Se exige que atienda la propia persona consultante y no otra, y se confirma que sea ella.
    \item \textbf{Verificar si tiene medidas.} A veces figura en la solicitud y otras no. Se lee lo que posea la persona consultante y se le pide que aporte la documentación correspondiente. Una vez asegurado que existen medidas y que están vigentes, sigue el proceso normal.
\end{enumerate}

% TODO: contenido incompleto o ambiguo en las notas originales (frase "le pedimos que pase fundamental": no queda claro qué se solicita exactamente)

\begin{important}
Si algo en este proceso falla, se debe llamar a los responsables de la comisión (a cualquiera de los cuatro) sin importar el día ni la hora, sea sábado, domingo o de noche. Aunque se lea la medida y no se la entienda, no se sepa para quién es o el nombre esté mal, igualmente corresponde llamar.
\end{important}

\section{Cuadro Cornell del capítulo}

\begin{longtable}{
    >{\raggedright\arraybackslash}p{0.26\textwidth}
    >{\raggedright\arraybackslash}p{\dimexpr\textwidth-0.26\textwidth-4\tabcolsep\relax}
}
\toprule
\textbf{Preguntas / palabras clave} & \textbf{Notas / respuestas} \\
\midrule
\endhead

\textbf{¿Cuándo se activa el protocolo de violencia?} &
Cuando llega un formulario que consigna una causa de violencia; se actúa de inmediato. \\[1.5em]

\textbf{¿Cómo se realiza la llamada a la persona consultante?} &
De manera cubierta, por una solicitud genérica de la UBA, sin dar muchos datos ni mencionar el patrocinio gratuito, y confirmando que atienda la propia persona. \\[1.5em]

\textbf{¿Qué se verifica después?} &
Si tiene medidas y si están vigentes, leyendo lo que posea; una vez asegurado, el proceso sigue normalmente. \\[1.5em]

\textbf{¿Qué se hace si algo falla?} &
Se llama a cualquiera de los cuatro responsables, a cualquier día y hora, aunque no se entienda la medida o haya errores. \\[1.5em]

\midrule
\multicolumn{2}{p{\dimexpr\textwidth-2\tabcolsep\relax}}{%
\textbf{Resumen:} Ante un formulario con causa de violencia se llama de inmediato y de forma cubierta a la persona consultante, se confirma su identidad y se verifica la vigencia de las medidas; ante cualquier falla se llama a los responsables sin importar el día ni la hora.} \\
\bottomrule
\end{longtable}

% ============================================================
\chapter[Orientación práctica: traslado]{Orientación práctica: traslado al exterior con un hijo}
% ============================================================

En una consulta se orientó a un progenitor que proyectaba salir del país con su hijo, existiendo acuerdo con el otro progenitor para firmar las autorizaciones.

\begin{itemize}
    \item Cuando se tenga una \textbf{fecha cierta} de salida (por ejemplo, la compra del pasaje), se debe volver a la comisión para formalizar el acuerdo por escrito.
    \item Debe documentarse tanto la salida como el \textbf{cambio de centro de vida} del niño, para evitar problemas posteriores, como una denuncia de que el progenitor se lo llevó.
    \item Conviene hacerlo con tiempo, porque si no los inconvenientes aparecen cuando la salida está próxima.
    \item Se advierte que enero incide en los plazos de tribunales, por lo que los trámites se estimaron para febrero o marzo, y que el trabajo comenzaría en diciembre.
\end{itemize}

% TODO: contenido incompleto o ambiguo en las notas originales (referencia a "7 tribunales igual 7 meses va a ser para febrero, marzo": la transcripción no es clara)

\end{document}
